%% pln-review: A Formal Review of Probabilistic Logic Networks
%%
%% Render with:
%%   lualatex -shell-escape pln-review.tex && bibtex pln-review && lualatex -shell-escape pln-review.tex && lualatex -shell-escape pln-review.tex

\documentclass[11pt,openany]{report}
\usepackage{pln-common}

\title{A Formal Review of\\
  \emph{Probabilistic Logic Networks}\\[4pt]
  Ben Goertzel, Matthew Ikl\'{e}, Izabela Freire Goertzel, Ari Heljakka (2009)\\[12pt]
  \large\texttt{(DRAFT --- \today)}}
\author{Zar and Oru\v{z}i\thanks{%
  Oru\v{z}i denotes AI-assisted collaboration using ChatGPT, Codex, and Claude Code.}}
\date{}

\begin{document}
\maketitle

\begin{abstract}
This document is a chapter-by-chapter formal review of the PLN
book~\cite{goertzel2009pln}.  For each chapter, we evaluate what formalizes
cleanly in Lean~4, what required correction, what fails outright
(counterexamples), and what remains heuristic or operational.

This review is \emph{book-relative}: it follows the PLN book's structure
and uses the book as the reference against which the formalization is
measured.  For the self-contained mathematical refoundation, see
\emph{The World-Model Calculus}~\cite{xiPLN2026}.

Every claim of ``formalized'' means: a Lean~4 theorem with zero
\code{sorry}, building against Mathlib.  Every claim of ``corrected''
means: the book's original statement is shown to be false in general,
and a side-conditioned replacement is proved.  Every claim of
``counterexample'' means: a concrete Lean witness that the stated property
fails.
\end{abstract}

\tableofcontents


%% ═══════════════════════════════════════════════════════════════════
\chapter{How to Read This Review}
\label{ch:methodology}
%% ═══════════════════════════════════════════════════════════════════

\section{Verdict Labels}

Every theorem, formula, or claim from the PLN book is assigned one of five
labels:

\begin{description}
\item[\verdictClean{}] The book's claim formalizes directly as stated, with
  at most minor notational changes.

\item[\verdictCorrected{}] The book's claim is incorrect as stated but
  admits a corrected formulation with explicit side conditions.  The
  original is false in general; the corrected version is proved.

\item[\verdictAssumption{}] The claim holds under assumptions that the book
  left implicit.  The formalization makes these assumptions explicit and
  proves the claim conditional on them.

\item[\verdictNoGo{}] A formal counterexample or no-go theorem showing
  that the book's claim fails in general.  These are constructive Lean
  witnesses.

\item[\verdictHeuristic{}] The claim is operational, heuristic, or
  otherwise not amenable to theorem-level formalization.  This includes
  runtime strategies, performance observations, and system-engineering
  advice.
\end{description}

\section{What Counts as ``Covered''}

A PLN book chapter is \emph{covered} when:
\begin{enumerate}
\item Its core mathematical claims have been formalized (possibly with
  corrections).
\item Its boundary conditions have been identified (counterexamples where
  claims fail).
\item Its operational/heuristic content has been explicitly marked as such.
\end{enumerate}

A chapter is \emph{not} covered merely by having been read or discussed.
Coverage requires theorem-level Lean artifacts.

\section{Relation to \xiPLN{}}

Full theorem statements and proofs live in the companion volume,
\emph{The World-Model Calculus}~\cite{xiPLN2026}.
This review provides:
\begin{itemize}
\item The book-relative context (what the original chapter was trying to do).
\item The evaluation (what succeeded, what failed, what changed).
\item Pointers to the corresponding \xiPLN{} chapters and Lean modules.
\end{itemize}


%% ═══════════════════════════════════════════════════════════════════
\chapter{Global Summary}
\label{ch:summary}
%% ═══════════════════════════════════════════════════════════════════

\begin{longtable}{p{0.08\textwidth}p{0.22\textwidth}p{0.22\textwidth}p{0.38\textwidth}}
\textbf{Ch.} & \textbf{Topic} & \textbf{Verdict} & \textbf{Key finding} \\
\hline
1 & Introduction & \verdictClean{} & Conceptual framing formalized as explicit semantic stack split. \\
2 & Knowledge repr. & \verdictClean{} & PLNQuery, WorldModel, WorldModelSigma capture the KR layer precisely. \\
3 & Experiential sem. & \verdictClean{} & Evidence-as-counts and revision algebra formalized exactly. \\
4 & Indefinite TVs & \verdictCorrected{} & Walley predictive and Bayesian credible intervals must be separated. \\
5 & FO ext.\ inference & \verdictCorrected{} & Classical formulas survive in corrected scope: BN exactness under screening-off/positivity, with no-go results showing that scope is necessary. \\
6 & FO ext.\ + ITV & \verdictClean{} & ITV transport through WM/OSLF bridge formalized. \\
7 & Distributional TVs & \verdictClean{} & Beta/Dirichlet layer, Kyburg flattening, de~Finetti bridge. \\
8 & Error magnification & \verdictCorrected{} & Threshold transport needs explicit boundary conditions. \\
9 & Large-scale inf. & \verdictNoGo{} & Multiple no-go results; corrected positives under explicit assumptions. \\
10 & Higher-order ext. & \verdictCorrected{} & Some unconditional equalities fail; side-conditioned replacements proved. \\
11 & Quantifiers + ITV & \verdictAssumption{} & Stable finite-domain bundle formalized; arbitrary-domain infinitary semantics already exist, but full parity is future work. \\
12 & Intensional inf. & \verdictCorrected{} & Settled in corrected typed three-channel form: ASSOC/PAT grounding, Solomonoff log-ratio bridge, selector endpoints, and explicit binary mixed-policy no-go. \\
13 & Inference control & \verdictClean{} & Selector specs, ranking, coverage formalized with canaries. \\
14 & Temporal/causal & \verdictClean{} & Event calculus grounding and temporal operators formalized. \\
App.\ A & PLN vs.\ NARS & \verdictClean{} & Consolidated correspondence package with confidence transforms. \\
\hline
\end{longtable}


%% ═══════════════════════════════════════════════════════════════════
%% CHAPTER REVIEWS
%% ═══════════════════════════════════════════════════════════════════

\chapter{Chapter 1: Introduction}
\label{ch:review-ch1}

\section{What the Book Is Trying to Do}

Chapter~1 introduces PLN's ambitions: a general-purpose uncertain reasoner
combining logical structure with probabilistic semantics.  It motivates
truth values (strength and confidence), evidence-based reasoning, and the
aspiration to handle deduction, induction, and abduction within a unified
framework.

\section{What Formalizes Cleanly}

The conceptual framing---separating probabilistic semantics from evidence
algebra from operational views---has been formalized as an explicit semantic
stack split:

\begin{itemize}
\item \textbf{Layer 1} (world-model posterior): \codepath{Mettapedia/Logic/PLNWorldModel.lean}
\item \textbf{Layer 2} (evidence algebra): \codepath{Mettapedia/Logic/EvidenceClass.lean},
  \codepath{Mettapedia/Logic/EvidenceQuantale.lean}
\item \textbf{Layer 3} (operational views): \codepath{Mettapedia/Logic/PLNIndefiniteTruth.lean}
\item \textbf{Decision tree}: \codepath{Mettapedia/Logic/SemanticsDecisionTree.lean}
\end{itemize}

\section{What Required Correction}

The book does not explicitly separate the three layers, leading to the
``semantic drift'' discussed in \xiPLN{} Chapter~1.  The correction is
not a change in formulas but a change in \emph{architecture}: making the
layer boundaries explicit.

\section{Counterexamples}

None for this chapter.  The introduction is interpretive, not
formula-bearing.

\section{What Remains Heuristic}

The book's motivational examples and historical context are expository
and not subject to formalization.

\section{Verdict}

\verdictClean{}.  The introduction's conceptual content formalizes well
once the three-layer separation is made explicit.  The refoundation is
architecturally cleaner but preserves the original intuitions.


\chapter{Chapter 2: Knowledge Representation}
\label{ch:review-ch2}

\section{What the Book Is Trying to Do}

Chapter~2 defines PLN's knowledge representation: atoms, links, truth
values, and the Atomspace data structure.  It establishes how uncertain
knowledge is stored and queried.

\section{What Formalizes Cleanly}

\begin{itemize}
\item \code{PLNQuery}: atoms, links, conditional links --- \codepath{Mettapedia/Logic/PLNWorldModel.lean}
\item \code{WorldModel}: the query interface --- same file
\item \code{WorldModelSigma}: side-condition-bundled models --- same file
\item Rewrite judgments: \codepath{Mettapedia/Logic/PLNWorldModelCalculus.lean}
\end{itemize}

\section{What Required Correction}

The book treats the Atomspace as a flat store of truth-valued links.
The formalization adds:
\begin{itemize}
\item Typed query language separating prop/link/linkCond.
\item Explicit side-condition contexts for structural assumptions.
\item World-model revision as a first-class operation.
\end{itemize}

\section{Counterexamples}

None.

\section{What Remains Heuristic}

The Atomspace as a software system (indexing, storage, concurrency) is
not formalized.  The formalization captures the mathematical interface.

\section{Verdict}

\verdictClean{}.  The KR layer formalizes directly with natural type-theoretic
representations.


\chapter{Chapter 3: Experiential Semantics}
\label{ch:review-ch3}

\section{What the Book Is Trying to Do}

Chapter~3 introduces the idea that PLN truth values are grounded in
\emph{evidence}: observations of positive and negative instances.
The core representation is the evidence pair $(n^+, n^-)$ where
$n^+$ counts positive observations and $n^-$ counts negative ones.
Revision---combining two bodies of evidence---is componentwise addition:
\[
  (n_1^+, n_1^-) \hplus (n_2^+, n_2^-) = (n_1^+ + n_2^+,\; n_1^- + n_2^-)
\]
The \emph{strength} is the ratio $s = n^+ / (n^+ + n^-)$ and the
\emph{confidence} is $c = n / (n + k)$ where $n = n^+ + n^-$.

The book's insight---that evidence counts are the right primitive, not
point probabilities---is exactly right.  De Finetti's theorem gives the
mathematical justification: for exchangeable binary sequences, the counts
are sufficient statistics.

\section{What Formalizes Cleanly}

\begin{itemize}
\item Evidence type $(n^+, n^-)$ with componentwise operations ---
  \codepath{Mettapedia/Logic/EvidenceClass.lean}
\item Revision as additive monoid (associative, commutative, unit $(0,0)$) ---
  same file
\item Information ordering (coordinatewise $\le$) --- same file
\item Complete world-model posterior (Dirichlet pseudo-counts over $2^n$ worlds) ---
  \codepath{Mettapedia/Logic/PLNJointEvidence.lean}
\item Evidence compositionality:
  $\mathrm{ev}(W_1 + W_2, q) = \mathrm{ev}(W_1, q) \hplus \mathrm{ev}(W_2, q)$ ---
  \codepath{Mettapedia/Logic/PLNWorldModel.lean}
\item Fr\'{e}chet bounds: $\max(0, s_A + s_B - 1) \le s_{A \wedge B} \le \min(s_A, s_B)$ ---
  \codepath{Mettapedia/Logic/PLNFrechetBounds.lean}
\end{itemize}

\section{What Required Correction}

The book is essentially correct here.  The formalization goes beyond
the book by discovering additional algebraic structure:
\begin{itemize}
\item \textbf{Quantale structure:} a tensor product $\otimes$ on evidence
  that distributes over joins, capturing compositional dependent-evidence
  combination.
\item \textbf{Heyting algebra:} a relative pseudo-complement
  $a \Rightarrow b$ satisfying $c \le (a \Rightarrow b)$ iff $c \otimes a \le b$.
  This means the evidence lattice supports intuitionistic logic.
\item \textbf{Totality gate:} $\Evidence$ has incomparable elements
  (e.g., $(3,1)$ and $(1,3)$), so no faithful order-embedding into $\mathbb{R}$
  exists.  A single real number cannot capture all evidence information.
\end{itemize}

These are \emph{additions}, not corrections.  The book planted the seed;
the formalization grew the full algebraic tree.

\section{Counterexamples}

\begin{itemize}
\item \code{threshold\_or\_counterexample}: because evidence values can be
  incomparable, disjunction threshold acceptance fails without assuming
  total-order comparability.
\end{itemize}

\section{What Remains Heuristic}

None.  This is the most naturally formal chapter in the book.

\section{Verdict}

\verdictClean{}.  The evidence algebra is the strongest part of the
original PLN book, and it formalizes beautifully.  The formalization
enriches it with quantale and Heyting structure that the book did not
explore but that the evidence representation naturally supports.


\chapter{Chapter 4: Indefinite Truth Values}
\label{ch:review-ch4}

\section{What the Book Is Trying to Do}

Chapter~4 introduces indefinite truth values (ITVs): interval-valued
representations of uncertainty that go beyond point estimates.

\section{What Formalizes Cleanly}

\begin{itemize}
\item ITV core with lower/upper/width/credibility coordinates ---
  \codepath{Mettapedia/Logic/PLNIndefiniteTruth.lean}
\item Walley IDM predictive interval constructors ---
  same file, explicit \code{walleyIDM} constructors
\item ITV hypercube transport ---
  \codepath{Mettapedia/Logic/PLNWorldModelITVHypercube.lean}
\end{itemize}

\section{What Required Correction}

The book uses the notation $[L, U]$ without always specifying which
interval semantics is intended.  The formalization \emph{insists} on
explicit separation:

\begin{itemize}
\item \textbf{Bayesian credible intervals} (posterior-dependent, prior-sensitive).
\item \textbf{Walley IDM predictive intervals} (prior-free, wider).
\item \textbf{Typed selector choice}: the user declares which semantics.
\end{itemize}

This is a \emph{strengthening}: the formalization is more precise than
the book, not a rejection of the book's ideas.

\section{Counterexamples}

\begin{itemize}
\item \code{canary\_ch11\_rule4\_not\_equivalent\_itvPath}: certain ITV
  transport paths are not equivalent, demonstrating that conflating
  interval semantics can produce different results.
\end{itemize}

\section{What Remains Heuristic}

The book's discussion of ``intuitively reasonable'' interval widths
is expository rather than mathematical.

\section{Verdict}

\verdictCorrected{}.  The ITV idea was right; the formalization
\emph{strengthens} it by separating semantics that the book conflated.


\chapter{Chapter 5: First-Order Extensional Inference}
\label{ch:review-ch5}

\section{What the Book Is Trying to Do}

Chapter~5 presents the core PLN inference rules---deduction, induction,
abduction, and revision---with strength formulas for computing truth values
of derived links.  The deduction formula, for instance, computes
$P(C \mid A)$ from $P(B \mid A)$, $P(C \mid B)$, $P(B)$, and $P(C)$:
\[
  s_{AC} = s_{AB} \cdot s_{BC} + (1 - s_{AB}) \cdot
  \frac{s_C - s_B \cdot s_{BC}}{1 - s_B}
\]
under the assumption that $C$ is conditionally independent of $A$
given $B$ and given $\lnot B$.

\section{What Formalizes Cleanly}

\begin{itemize}
\item Deduction strength formula: \code{plnDeductionStrength} ---
  \codepath{Mettapedia/Logic/PLNDeduction.lean}
\item Induction and abduction strength formulas ---
  \codepath{Mettapedia/Logic/PLNDerivation.lean}
\item Full inference rule set ---
  \codepath{Mettapedia/Logic/PLNInferenceRules.lean}
\item MeTTa truth function parity ---
  \codepath{Mettapedia/Logic/PLNMettaTruthFunctions.lean}
\item BN chain exactness: \code{chainBN\_plnDeductionStrength\_exact} ---
  \codepath{Mettapedia/Logic/PLNBayesNetFastRules.lean}
\item BN fork exactness: \code{forkBN\_plnDeductionStrength\_exact} ---
  same file
\end{itemize}

The deduction formula is proved \emph{exact} for the chain BN
$A \to B \to C$ and the fork BN $A \leftarrow B \to C$.  That is:
under the respective d-separation conditions, the formula above equals
$P(C \mid A)$ exactly---it is not an approximation.

\section{What Required Correction}

The book \emph{does} mention conditional independence as a motivation
for the deduction formula (Section~5.1 discusses the total probability
decomposition, which implicitly assumes screening-off).  However, the
precise conditions under which each rule is exact are not always stated
as formal hypotheses.  The formalization's contribution is to recover
the correct theoremic scope of the classical formulas by making these
conditions fully explicit:

\begin{itemize}
\item Each rule carries a $\Sigma$-context stating the required
  conditional independences (e.g., $C \perp\!\!\!\perp A \mid B$ and
  $C \perp\!\!\!\perp A \mid \lnot B$ for deduction).
\item Positivity conditions ($s_B > 0$, $s_B < 1$) are stated as
  explicit hypotheses rather than tacit assumptions.
\item In BN model classes, the needed side conditions are discharged by
  d-separation/local-Markov facts, so the rules are proved as
  \emph{derived theorems} rather than treated as axioms or heuristics.
\end{itemize}

This is a corrected statement rather than an assumption patch.  The
classical extensional formulas survive, but only in the precise scope
determined by screening-off, d-separation, and positivity, with
exactness recovered on the corresponding BN-derived envelopes.

\section{Counterexamples}

\begin{itemize}
\item The local-completeness no-go theorem (\xiPLN{} Chapter~6) shows
  that no function of five local evidence values computes exact link
  evidence for \emph{arbitrary} world models.  This establishes that
  the screening-off conditions are not merely convenient---they are
  \emph{necessary} for correctness.
\end{itemize}

This is different in kind from Chapter~11's current finite-domain front
door: Chapter~5's side conditions are the mathematical content of the
exact theorem, not a packaging gap between the reviewed chapter and a
deeper semantic layer.

\section{What Remains Heuristic}

The choice of \emph{when} to apply each rule (which links to try
deduction on, in what order) is inference control, covered in
Chapter~13.

\section{Verdict}

\verdictCorrected{}.  The book-era rule statements were too broad, but
the classical extensional formulas survive in corrected form, with BN
exactness theorems under explicit screening-off / d-separation /
positivity conditions and no-go results showing that this scope is
mathematically necessary.


\chapter{Chapter 6: First-Order Extensional Inference with ITVs}
\label{ch:review-ch6}

\section{What the Book Is Trying to Do}

Chapter~6 extends the first-order rules to operate on indefinite truth
values (intervals) rather than simple truth values.

\section{What Formalizes Cleanly}

\begin{itemize}
\item ITV semantics transported through WM rewrites ---
  \codepath{Mettapedia/Logic/PLNWMOSLFBridgeITV.lean}
\item One-call API bundles for ITV-coordinated inference ---
  \codepath{Mettapedia/Logic/PLNCanonicalAPI.lean}
\item Coordinate-specialized coherence theorems ---
  same file
\end{itemize}

\section{What Required Correction}

No fundamental correction needed.  The formalization adds:
\begin{itemize}
\item Typed OSLF bridge with explicit selector specialization.
\item Coordinate-by-coordinate transport
  (lower/upper/credibility/width/strength).
\end{itemize}

\section{Counterexamples}

None for this chapter specifically.

\section{What Remains Heuristic}

None.

\section{Verdict}

\verdictClean{}.  The ITV extension of first-order rules formalizes
via the typed bridge infrastructure.


\chapter{Chapter 7: Distributional Truth Values}
\label{ch:review-ch7}

\section{What the Book Is Trying to Do}

Chapter~7 presents distributional truth values: truth values as entire
distributions (Beta, Dirichlet) rather than point estimates or intervals.
It introduces the Kyburg flattening operation.

\section{What Formalizes Cleanly}

\begin{itemize}
\item Context-aware strength as Beta posterior mean:
  \code{strengthWith\_eq\_beta\_posterior\_meanENN} ---
  \codepath{Mettapedia/Logic/HigherOrder/PLNKyburgReduction.lean}
\item Evidence aggregation as conjugate update:
  \code{hplus\_is\_bayesian\_update} --- same file
\item Kyburg flattening: \code{kyburg\_flattening},
  \code{expectation\_consistency}, \code{kyburg\_no\_advantage} ---
  same file and \codepath{Mettapedia/ProbabilityTheory/HigherOrderProbability/KyburgFlattening.lean}
\item Monad associativity for flattening:
  \code{flatten\_associativity}, \code{flatten\_associativity\_kernel} ---
  \codepath{Mettapedia/ProbabilityTheory/HigherOrderProbability/GiryMonad.lean}
\item De~Finetti singleton bridge into Kyburg flattening ---
  \codepath{Mettapedia/Logic/HigherOrder/PLNKyburgReduction.lean}
\item Worked example (uniform prior, evidence 3:1) ---
  same file
\end{itemize}

\subsection{Normal-Gamma Extension for Continuous Domains}

The distributional story extends beyond Beta/Dirichlet to continuous
observations via Normal-Gamma conjugate priors.  Evidence becomes the
sufficient statistic $(n, \sum x_i, \sum x_i^2)$; revision is again
componentwise addition.  The key compositionality theorem---batch
replay equals sequential Bayesian update---is proven:

\begin{quote}
\code{posterior\_hplus\_of\_realizable}: \\
$\text{posterior}(\pi, e_1 \hplus e_2)
  = \text{posterior}(\text{posterior}(\pi, e_1), e_2)$
\end{quote}

See \codepath{Mettapedia/Logic/EvidenceNormalGamma.lean}.

Six applied demos (broken sensor, widget factory, Kalman filter, gas
sensor drift, steel faults, NAB anomaly) exercise these properties
end-to-end on synthetic and real UCI/NAB data, totaling $\sim$2{,}500
lines of Lean with 0~sorry; see the WM-PLN book, Ch.~7.  A seventh
script-level case study extends the same update semantics to the public
Good Judgment Project / ACE forecasting stream: reliability-weighted
aggregation improves substantially over mean/median pooling on held-out
year-4 questions, and it also improves Brier/log score over stronger
top-$k$ log/logit pooling baselines while matching their accuracy.
Empirically, the role split should be read as two flagships plus two
supporting checks: GJP forecasting and Gas drift governance are the main
applied stories, Steel is a supporting same-model-class abstention
validation, and NAB is an honest negative boundary / case study.
An explicit extremized-mean comparator is also weaker than the
top-$k$ logit family, so the WM gain is not just an artifact of beating
plain averaging.
The same WM-weighted pools also improve selective prediction AURC over
those stronger baselines, while WM-on-top-$k$ hybrids do not improve
further, which is evidence that the gain comes from continuous
reliability weighting rather than hard forecaster pruning.  Nightly
batch consolidation matches online updates to machine precision, and
scoped forgetting of older years continues to improve recalibration even
for the stronger pooling family.  A small support-layer result also
appears: on multi-option year-4 questions, gating the strong top-$k$
log pool by WM agreement halves AURC from 0.001095 to 0.000541 while
preserving accuracy, whereas the binary slice is an honest null because
the strong top-$k$ and WM-weighted pools already coincide everywhere.
That forecasting study is not a new
Lean theorem block, but it is a clean applied demonstration of the same
conjugate evidence and revision laws.

\section{What Required Correction}

The Kyburg ``no advantage'' result is a sharpening: the book presents
distributional truth values as potentially more powerful than simple
truth values, but the formalization shows that for expectation-based
decisions, they are equivalent (the flattening preserves expectation).

\section{Counterexamples}

None.

\section{What Remains Heuristic}

The book's discussion of when distributional truth values are
``worth the computational cost'' is a practical consideration,
not a theorem.

\section{Verdict}

\verdictClean{}.  This is one of the strongest formalizations.
The Kyburg flattening, de~Finetti bridge, conjugate-update story,
and Normal-Gamma continuous extension all formalize beautifully.


\chapter{Chapter 8: Error Magnification}
\label{ch:review-ch8}

\section{What the Book Is Trying to Do}

Chapter~8 discusses how errors accumulate through chains of inference
(``error magnification'') and proposes diagnostic approaches.

\section{What Formalizes Cleanly}

\begin{itemize}
\item View transport under WM query equivalence: strength, confidence,
  interpretation views all preserved ---
  \codepath{Mettapedia/Logic/PLNWorldModelCalculus.lean}
\item Confidence-threshold atom semantics and rewrite soundness ---
  \codepath{Mettapedia/Logic/PLNErrorMagnificationGrounding.lean}
\item One-call WM rewrite $\to$ OSLF truth transport $\to$ threshold
  acceptance --- \codepath{Mettapedia/Logic/PLNCanonicalAPI.lean}
\item WM consequence-rule layer with Kripke deontic lifts ---
  \codepath{Mettapedia/Logic/PLNWorldModelKripkeConsequence.lean}
\item Categorical WM transport (institution, hyperdoctrine,
  Beck-Chevalley) ---
  \codepath{Mettapedia/Logic/PLNWorldModelCategoricalBridge.lean}
\item Kripke and neighborhood proof-theoretic closure (K, KT, KD, K4,
  E, EMN, EMT, ED) ---
  \codepath{Mettapedia/Logic/PLNWorldModelKripkeCompleteness.lean},
  \codepath{Mettapedia/Logic/PLNWorldModelNeighborhoodCompleteness.lean}
\end{itemize}

\section{What Required Correction}

The book's error magnification discussion is largely qualitative
(graphs, chaos-plot intuitions).  The formalization replaces this with:

\begin{itemize}
\item Precise threshold transport theorems with explicit boundary conditions.
\item Double-damping strict gap characterization.
\item Explicit counterexamples for disjunction and diamond threshold
  transport.
\end{itemize}

\section{Counterexamples}

\begin{itemize}
\item \code{threshold\_or\_counterexample}: disjunction threshold fails
  without total-order comparability.
\item \code{threshold\_dia\_fails}: diamond threshold requires finite
  branching.
\item \code{double\_damping\_strict\_lt}: composed thresholds are
  strictly tighter than individual thresholds.
\end{itemize}

\section{What Remains Heuristic}

The book's chaos plots, sensitivity diagrams, and qualitative
error-propagation discussions are expository.  The formalized core is
the threshold transport discipline with its precise boundaries.

\section{Verdict}

\verdictCorrected{}.  The book's qualitative discussion is replaced by
precise transport theorems with explicit boundary conditions.  The
formalization is significantly more precise than the original.


\chapter{Chapter 9: Large-Scale Inference Strategies}
\label{ch:review-ch9}

This is the most complex chapter to review.  We break it into subsections
corresponding to the book's major topics.

\section{What the Book Is Trying to Do}

Chapter~9 covers strategies for scaling PLN to large knowledge bases:
batch inference, multideduction, BN augmentation, trail-free protocols,
and evidence pooling.

\subsection{9.2: Batch Inference / BRGI}

\paragraph{What formalizes.}
\begin{itemize}
\item Finite BRGI object with policy-sensitive optimum characterization:
  \code{exists\_optimal\_feasible\_set\_with\_policy\_characterization} ---
  \codepath{Mettapedia/Logic/PremiseSelectionBRGI.lean}
\item Graph-level BRGI with optimality-preserving refinement ---
  same file
\end{itemize}

\paragraph{Verdict.}
\verdictClean{} for the optimization-theoretic core.
\verdictHeuristic{} for the runtime implementation.

\subsection{9.3: Multideduction}

The book proposes combining results from multiple deduction chains using
inclusion-exclusion.  The first two terms of IE give:
\[
  |A \cup B| \approx |A| + |B| - |A \cap B|
\]
But the question is: can $|A \cup B|$ be determined from these two-term
statistics alone?

\paragraph{What formalizes.}
\begin{itemize}
\item \textbf{Identifiability no-go:}
  \code{same\_first\_two\_terms\_different\_union} constructs two
  probability distributions agreeing on all first-two-term statistics
  but differing on $|A \cup B|$.  Therefore two-term IE is
  \emph{not identifiable}. ---
  \codepath{Mettapedia/Logic/PLNInclusionExclusionIdentifiability.lean}
\item \textbf{No universal predictor:}
  \code{no\_universal\_union\_predictor\_from\_first\_two\_terms} ---
  same file.  No additive correction term works for all models.
\item \textbf{Residual decomposition:}
  $|A \cup B| = T_2(A,B) + R(A,B)$ where $T_2$ is the two-term estimate
  and $R$ is the residual.  If $R$ is known or bounded, the corrected
  estimate is exact:
  \code{correctedEstimate\_agrees\_of\_residualAssumption} ---
  \codepath{Mettapedia/Logic/PLNMultideductionResidual.lean}
\item \textbf{N-way generalization:} the residual decomposition extends
  to $n$-way unions. ---  same file
\end{itemize}

\paragraph{Verdict.}
\verdictNoGo{} for two-term IE as an identifiable estimator.
\verdictAssumption{} for the corrected residual-based approach (requires
knowing or bounding the residual $R$).

\subsection{9.4: BN Augmentation}

\paragraph{What formalizes.}
\begin{itemize}
\item Class-packaged BN d-separation discharge ---
  \codepath{Mettapedia/Logic/PLNBNLocalMarkovPackages.lean}
\item One-call selector $\to$ rewrite $\to$ threshold composition ---
  \codepath{Mettapedia/Logic/PLNCanonicalAPI.lean}
\item Chain/fork/collider concrete endpoints ---
  \codepath{Mettapedia/Logic/PLNXiDerivedBNRules.lean}
\end{itemize}

\paragraph{Verdict.}
\verdictClean{}.  One of the strongest formalizations: exact BN query
answering with class-packaged assumption discharge.

\subsection{9.5: Trail-Free / OBI Protocols}

The book proposes trail-free inference: updating truth values without
tracking which inference paths contributed, to avoid exponential trail
bookkeeping.  The hope is that repeated application converges to a
fixed point.

\paragraph{What formalizes.}
\begin{itemize}
\item \textbf{Non-stabilization counterexample:}
  \code{orbit\_not\_eventually\_constant} constructs a two-state system
  where the undamped trail-free update enters a 2-cycle
  $s_0 \to s_1 \to s_0 \to \cdots$ and never reaches a fixed point.
  The construction is explicit and decidable. ---
  \codepath{Mettapedia/Logic/PLNTrailFreeDynamicsCounterexample.lean}
\item \textbf{Damped convergence:} with a damping factor
  $\alpha \in (0,1)$ and eventual fresh evidence, the orbit
  eventually stabilizes:
  \code{damped\_sp\_spn\_eventually\_constant\_of\_eventual\_fresh}.
  The inconsistency measure drops to zero. ---
  \codepath{Mettapedia/Logic/PLNTrailFreeDampedConvergence.lean}
\item \textbf{Bounded-gap fairness:} if fresh evidence arrives with
  bounded gaps (at most $g$ steps between arrivals), stabilization
  occurs within an explicit bound:
  \code{inconsistency\_zero\_within\_bound\_of\_boundedGapFresh}. ---
  same file
\item \textbf{No-fresh negative:} without fresh evidence, even the
  damped protocol may not stabilize:
  \code{dampedOrbit\_noFresh\_not\_eventually\_constant}. ---
  same file
\end{itemize}

\paragraph{Verdict.}
\verdictNoGo{} for undamped trail-free protocols.
\verdictAssumption{} for damped variants (requires fresh evidence
with bounded gaps).

\subsection{Evidence Pooling}

\paragraph{What formalizes.}
\begin{itemize}
\item Pooling axioms non-uniqueness:
  \code{maxPoolingOperator\_ne\_fuse} ---
  \codepath{Mettapedia/Logic/PremiseSelectionExternalBayesianity.lean}
\item Corrected uniqueness: external Bayes + total additivity ---
  same file
\end{itemize}

\paragraph{Verdict.}
\verdictNoGo{} for pooling axioms alone.
\verdictAssumption{} with external Bayesianity.

\subsection{Stochastic Experiment Layer}

\paragraph{What formalizes.}
\begin{itemize}
\item Blackwell factorization and expected utility ---
  \codepath{Mettapedia/Logic/PLNWorldModelExperimentStochastic.lean}
\item Weighted source policies and trusted-gate transfer ---
  same file
\end{itemize}

\paragraph{Verdict.}
\verdictClean{}.

\section{Overall Chapter 9 Verdict}

\verdictNoGo{} for several core claims (identifiability, trail-free
stabilization, pooling uniqueness).  \verdictAssumption{} for
corrected positive replacements.  \verdictClean{} for the BN
augmentation track.  Operational runtime implementations remain
\verdictHeuristic{}.


\chapter{Chapter 10: Higher-Order Extensional Inference}
\label{ch:review-ch10}

\section{What the Book Is Trying to Do}

Chapter~10 extends PLN to reason about relationships between predicates
and concepts (higher-order extensional inference), including satisfying
sets and reductions from higher-order to first-order.

\section{What Formalizes Cleanly}

\begin{itemize}
\item HOI-to-FOI reduction via satisfying sets: for predicate $P$, the
  satisfying set $\mathrm{Sat}(P) = \{x \mid P(x)\}$ bridges
  higher-order and first-order ---
  \codepath{Mettapedia/Logic/HigherOrder/HigherOrderReduction.lean}
\item Weakness-based inheritance:
  $\mathrm{Inheritance}(A, B, \mu) = \mathrm{weakness}(A \cap B) / \mathrm{weakness}(A)$
  ---
  \codepath{Mettapedia/Logic/HigherOrder/Basic.lean}
\item Pointwise implication collapse: when $\forall x, R_1(A,x) \Rightarrow R_2(B,x)$
  holds, inheritance equals~$1$ ---
  same file
\item Perfect inheritance corner ($\mathrm{pos} \ne 0$, $\mathrm{neg} = 0
  \Rightarrow \mathrm{Inheritance} = \mathrm{pTrue}$) ---
  same file
\end{itemize}

\section{What Required Correction}

The book's higher-order composition rule
$\mathrm{Inh}(A,C) = f(\mathrm{Inh}(A,B), \mathrm{Inh}(B,C))$ is
stated unconditionally.  The formalization shows:
\begin{itemize}
\item The rule requires screening-off: $P(C \mid A, B) = P(C \mid B)$.
\item Positivity conditions ($P(B) \in (0,1)$) must hold.
\item The corrected version carries these as explicit hypotheses, matching
  the structure of the first-order deduction correction.
\end{itemize}

\section{Counterexamples}

Where the book's unconditional equalities fail, the formalization provides
concrete witnesses with specific evidence values showing the failure mode.

\section{What Remains Heuristic}

None of the core mathematical content.

\section{Verdict}

\verdictCorrected{}.  The reduction is sound under stated conditions;
several book-era unconditional claims are shown to need screening-off
and positivity side conditions, paralleling the first-order story.


\chapter{Chapter 11: Quantifiers with Indefinite Truth Values}
\label{ch:review-ch11}

\section{What the Book Is Trying to Do}

Chapter~11 extends PLN to quantified reasoning (``most birds fly,''
``few cats swim'') using indefinite truth values for quantifier
strength.

\section{What Formalizes Cleanly}

\begin{itemize}
\item Finite-model quantifier semantics with weakness connection ---
  \codepath{Mettapedia/Logic/PLNFirstOrder/QuantifierSemantics.lean}
\item Soundness theorems ---
  \codepath{Mettapedia/Logic/PLNFirstOrder/Soundness.lean}
\item Fuzzy quantifier semantics (QFM) ---
  \codepath{Mettapedia/Logic/PLNFirstOrder/FuzzyQuantifierSemantics.lean}
\item Fuzzy-to-ITV bridge ---
  \codepath{Mettapedia/Logic/PLNFirstOrder/FuzzyITVBridge.lean}
\item End-to-end canaries:
  \code{canary\_ch11\_rule\_family\_end\_to\_end\_ext},
  \code{fuzzyIntervalHolds\_strength\_of\_lower\_upper}
\item Regression target:
  \codepath{Mettapedia/Logic/PLNFirstOrder/QuantifierRegression.lean}
\end{itemize}

\section{What Required Correction}

No fundamental corrections.  The main clarification is scope: the
reviewed chapter bundle is finite-domain, but an arbitrary-domain
infinitary semantic layer already exists in Lean.

\section{Counterexamples}

\begin{itemize}
\item \code{canary\_ch11\_rule4\_not\_equivalent\_itvPath}: certain ITV
  transport paths are not equivalent, demonstrating that some book-era
  assumptions about ITV composition are too strong.
\end{itemize}

\section{What Remains Heuristic}

Arbitrary-domain infinitary semantics already exist in
\codepath{Mettapedia/Logic/PLNFirstOrder/Infinite.lean}; what remains
beyond the current reviewed scope is stronger metatheory, chapter-level
canary/regression parity, and measure-theoretic strengthening for
probabilistic infinite spaces.

\section{Verdict}

\verdictAssumption{}.  Solid finite-domain theorem bundle, plus a real
arbitrary-domain infinitary semantic layer that is not yet promoted to
full chapter-level parity.


\chapter{Chapter 12: Intensional Inference}
\label{ch:review-ch12}

\section{What the Book Is Trying to Do}

Chapter~12 presents intensional inheritance: reasoning based on shared
properties (intension) rather than shared instances (extension).
It defines ASSOC and PAT channels and proposes mixing them with a
parameter.  The formal result is not verbatim preservation of that
binary presentation, but a corrected typed three-channel account with
an explicit no-go against collapsing back to the old binary law.

\section{What Formalizes Cleanly}

\begin{itemize}
\item Typed WM grounding for intensional inheritance ---
  \codepath{Mettapedia/Logic/PLNIntensionalWorldModel.lean}
\item Solomonoff-linked semantic transport ---
  \codepath{Mettapedia/Logic/IntensionalInheritanceSolomonoffBridge.lean}
\item Canonical \code{AssocPatSemanticModel} packaging ---
  \codepath{Mettapedia/Logic/PLNCanonicalAPI.lean}
\item Selector-specialized end-to-end endpoints ---
  same file
\item Canary and regression targets ---
  \codepath{Mettapedia/Logic/PLNIntensionalCanary.lean},
  \codepath{Mettapedia/Logic/PLNIntensionalRegression.lean}
\end{itemize}

\section{What Required Correction}

The book proposes mixing extensional and intensional inheritance with a
single parameter~$\alpha$:
\[
  \mathrm{InhInt}(A, B) = \alpha \cdot \mathrm{Ext}(A, B)
  + (1 - \alpha) \cdot \mathrm{IntAssoc}(A, B)
\]

The formalization shows this binary mixture is explicitly false as a
general law in the corrected typed setting:

\begin{itemize}
\item The ASSOC channel uses mutual-information semantics:
  $P_{\mathrm{ASSOC}}(W \mid F) = P(W) \cdot 2^{I(F;W)}$.
\item The PAT channel captures structural pattern proximity that mutual
  information misses.
\item The concrete witness $W_{\mathrm{demo}}$ with scores $(3, 2, 1)$
  for (ext, ASSOC, PAT) channels shows that the PAT channel contributes
  nontrivially and the mixed result cannot be expressed as any function
  of only the extensional and ASSOC values.
\end{itemize}

The corrected approach is not a patched binary rule.  It is a typed
three-channel model (extensional, ASSOC, PAT) with separate evidence
extraction per channel, a semantic model that packages all three, and a
constructive no-go theorem retiring the old binary law.

\section{Counterexamples}

\begin{itemize}
\item \code{binary\_mixed\_policy\_collapse\_no\_go}:
  the old binary mixed law is explicitly false in the corrected typed
  setting.  The constructive witness is
  \code{no\_binary\_mixed\_policy\_for\_assocPat\_fixture}, where
  $\mathrm{pat\_channel\_nontrivial}$ and
  $\mathrm{mixed\_not\_assoc\_only}$ are both proved.
\end{itemize}

\section{Claim-by-Claim Audit}

\begin{itemize}
\item \textbf{$\mathrm{InheritanceSort} = \mathrm{extensional} \mid
  \mathrm{intensionalAssoc} \mid \mathrm{intensionalPAT} \mid \mathrm{mixed}$,
  with separate channel projections.} \textsc{Proved}.  Endpoints:
  \code{InheritanceSort}, \code{InheritanceQueryFamily},
  \code{extensionalEvidence}, \code{intensionalAssocEvidence},
  \code{intensionalPATEvidence}, \code{mixedEvidence}.
\item \textbf{$P_{\mathrm{ASSOC}}(W \mid F) = P(W) \cdot 2^{I(F;W)}$.}
  \textsc{Strengthened}.  The abstract information-theoretic formula is
  available as \code{goertzel\_formula}; the corrected typed WM layer
  exports it through \code{AssocScoreCorrespondence} and
  \code{assocEvidence\_eq\_scoreToEvidence\_of\_assocScore\_eq}.
\item \textbf{PAT is a complementary structural channel, not reducible to
  ASSOC syntax alone.} \textsc{Strengthened}.  Endpoints:
  \code{PATScoreCorrespondence},
  \code{patEvidence\_eq\_scoreToEvidence\_of\_patScore\_eq}, and
  \code{AssocPatSemanticModel}.
\item \textbf{Universal-mixture intensional inheritance is the
  base-2 log ratio of conditionals.} \textsc{Proved}.  Endpoints:
  \code{intensionalFromConditional\_eq\_log2\_ratio},
  \code{intensionalFromXiSemimeasure\_eq\_log2\_ratio}, and
  \code{intensionalFromXiGeom\_eq\_log2\_ratio}.
\item \textbf{The chapter's Kolmogorov / normalized-information-distance gloss.}
  \textsc{Strengthened}.  The exact formal endpoint is the
  universal-mixture log-ratio theorem above; the complexity-distance
  language is retained only as interpretation, not as the primary
  theorem statement.
\item \textbf{$\mathrm{InhInt}(A, B) = \alpha \cdot \mathrm{Ext}(A, B) +
  (1 - \alpha) \cdot \mathrm{IntAssoc}(A, B)$.} \textsc{Disproved}.
  Endpoints: \code{binary\_mixed\_policy\_collapse\_no\_go},
  \code{mixed\_not\_assoc\_only}, and
  \code{pat\_channel\_nontrivial}.
\item \textbf{Concrete mixed projections and their divergence are explicit.}
  \textsc{Proved}.  Endpoints:
  \code{canary\_ch12\_mixed\_extensional\_projection},
  \code{canary\_ch12\_mixed\_assoc\_projection}, and
  \code{canary\_ch12\_mixed\_projection\_non\_equivalent}.
\item \textbf{Typed ASSOC/PAT semantic packaging has selector-specialized
  end-to-end endpoints.} \textsc{Proved}.  Endpoints:
  \code{mixedPolicy\_assocPat},
  \code{end\_to\_end\_assocPat\_bayesNormal\_lower},
  \code{end\_to\_end\_assocPat\_bayesExact\_lower}, and
  \code{end\_to\_end\_assocPat\_walley\_lower}.
\end{itemize}

\section{Settlement}

No chapter-level formula above remains open: each is now explicitly
marked as proved, strengthened to the corrected typed statement, or
ruled out.  Any further alias polish would be presentational only and
would not change the mathematical status of the chapter.

\section{Verdict}

\verdictCorrected{}.  The book's binary extensional+ASSOC mixture is too
weak as stated.  In the corrected typed setting, ASSOC and PAT live as
separate channels with explicit semantic packaging, Solomonoff linkage,
selector-specialized endpoints, and a constructive no-go against
collapsing them to a binary mixed policy.


\chapter{Chapter 13: Aspects of Inference Control}
\label{ch:review-ch13}

\section{What the Book Is Trying to Do}

Chapter~13 discusses how to control PLN inference: which rules to apply,
which premises to select, how to allocate computational resources.

\section{What Formalizes Cleanly}

\begin{itemize}
\item Formal selector specifications ---
  \codepath{Mettapedia/Logic/PremiseSelectionSelectorSpec.lean}
\item Ranking transfer (optimality under distribution shift) ---
  \codepath{Mettapedia/Logic/PremiseSelectionOptimality.lean}
\item Ranking stability ---
  \codepath{Mettapedia/Logic/PremiseSelectionRankingStability.lean}
\item Coverage/submodularity with greedy guarantee ---
  \codepath{Mettapedia/Logic/PremiseSelectionCoverage.lean}
\item Inference-control core composition ---
  \codepath{Mettapedia/Logic/PLNInferenceControlCore.lean}
\item Canary and regression targets ---
  \codepath{Mettapedia/Logic/PLNInferenceControlCanary.lean},
  \codepath{Mettapedia/Logic/PLNInferenceControlRegression.lean}
\end{itemize}

\section{What Required Correction}

No fundamental corrections.  The formalization adds:
\begin{itemize}
\item Submodularity (not in the book) with the $(1-1/e)$ greedy
  approximation guarantee.
\item Coverage surrogate no-go: coverage and cardinality are insufficient
  for risk-sensitive selection.
\end{itemize}

\section{Counterexamples}

\begin{itemize}
\item \code{sameCoverageAndCard\_differentFalsePositiveRisk}: coverage
  surrogate limitation.
\end{itemize}

\section{What Remains Heuristic}

The systems-engineering aspects of inference control (scheduling,
resource allocation, attention) are operational and not theorem-level.

\section{Verdict}

\verdictClean{}.  The mathematical core of inference control formalizes
well, with the submodularity story going beyond the book.


\chapter{Chapter 14: Temporal and Causal Inference}
\label{ch:review-ch14}

\section{What the Book Is Trying to Do}

Chapter~14 extends PLN to temporal reasoning (before/after/during) and
causal inference (cause-effect relationships).

\section{What Formalizes Cleanly}

\begin{itemize}
\item Temporal predicate type: $\mathrm{TemporalPred}(D, T) = D \to T \to \mathrm{Prop}$
\item Simultaneous operators: $\mathrm{SimAND}(P, Q)(x, t) = P(x,t) \wedge Q(x,t)$
\item Sequential operators:
  $\mathrm{SeqAND}(P, Q)(x, t) = P(x, t) \wedge \exists \Delta > 0.\; Q(x, t + \Delta)$
\item Predictive implication:
  $\mathrm{PredImpl}(P, Q)(x, t) = P(x, t) \Rightarrow \exists \Delta.\; Q(x, t + \Delta)$
\item 3-node predictive chain with temporal bounds ---
  \codepath{Mettapedia/Logic/PLNTemporalCausalInference.lean}
\item Probabilistic event calculus ($\mathrm{Happens}$, $\mathrm{Initiates}$,
  $\mathrm{Terminates}$, $\mathrm{HoldsAt}$) ---
  \codepath{Mettapedia/Logic/PLNProbabilisticEventCalculus.lean}
\item One-call event endpoints ---
  \codepath{Mettapedia/Logic/PLNCanonicalAPI.lean}
\end{itemize}

\section{What Required Correction}

No corrections needed.  The formalization treats temporal/causal
inference as another world-model instance rather than a separate
ontology.  The key insight: temporal structure lives in the
world-model state, temporal operators live in the query language,
and the evidence/revision machinery operates unchanged.

\section{Counterexamples}

None.

\section{What Remains Heuristic}

The philosophical aspects of causation (interventionist vs.\
counterfactual vs.\ probabilistic) are not resolved by formalization;
the formalization provides the probabilistic-event-calculus substrate.
Granger-style temporal precedence is used to distinguish causation from
correlation within the world-model framework, but this is a pragmatic
rather than metaphysical choice.

\section{Verdict}

\verdictClean{}.  Compact core that formalizes as a WM instance.  The
temporal/causal layer is one of the cleanest demonstrations of the
world-model approach: a single kernel, many instantiations.


\chapter{Appendix A: Comparison of PLN and NARS}
\label{ch:review-appA}

\section{What the Book Is Trying to Do}

Appendix~A compares PLN's inference rules with NARS (Non-Axiomatic
Reasoning System), showing correspondences and differences.

\section{What Formalizes Cleanly}

\begin{itemize}
\item Consolidated PLN$\leftrightarrow$NARS comparison package ---
  \codepath{Mettapedia/Logic/PLNNARSRuleCorrespondence.lean}
\item Confidence transforms between PLN and NARS evidence ---
  \codepath{Mettapedia/Logic/NARSEvidenceBridge.lean}
\item Galois connection between evidence algebras ---
  \codepath{Mettapedia/Logic/NARSPLNGaloisConnection.lean}
\end{itemize}

\section{What Required Correction}

None.  The comparison is straightforward once both systems are
formalized.

\section{Counterexamples}

None.

\section{What Remains Heuristic}

The qualitative comparison of ``PLN philosophy vs.\ NARS philosophy''
is not theorem-level.

\section{Verdict}

\verdictClean{}.  The rule correspondence and evidence bridge are
fully formalized.


%% ═══════════════════════════════════════════════════════════════════
\appendix
%% ═══════════════════════════════════════════════════════════════════

\chapter{Full Theorem Ledger}
\label{app:ledger}

The complete theorem ledger is maintained in the Lean codebase.
Key entry points:

\begin{itemize}
\item \codepath{Mettapedia/Logic/PLNCore.lean} --- umbrella module
\item \codepath{Mettapedia/Logic/PLNCanonicalAPI.lean} --- canonical API (3806 lines)
\item \codepath{Mettapedia/Logic/README.md} --- status table and theorem index
\item \codepath{Mettapedia/Implementation/PLNParityChecklist.lean} --- MeTTa parity
\end{itemize}

Chapter-specific regression scripts:
\begin{itemize}
\item \codepath{scripts/check\_ch8\_neighborhood.sh}
\item \codepath{scripts/check\_ch9\_positive.sh}
\item \codepath{scripts/check\_ch11\_quantifiers.sh}
\item \codepath{scripts/check\_ch11\_fuzzy\_syllogism.sh}
\item \codepath{scripts/check\_ch12\_intensional.sh}
\item \codepath{scripts/check\_ch13\_inference\_control.sh}
\end{itemize}


\chapter{Corrected Formulas}
\label{app:corrected}

\begin{longtable}{p{0.30\textwidth}p{0.30\textwidth}p{0.30\textwidth}}
\textbf{Book formula} & \textbf{Issue} & \textbf{Corrected version} \\
\hline
Deduction: $s_{AC} = f(s_{AB}, s_{BC}, s_B, s_C)$ unconditionally &
Missing screening-off / positivity scope &
Same formula, but only under explicit screening-off and positivity side conditions \\
\hline
ITV: $[L, U]$ (unspecified semantics) &
Conflates Bayesian credible and Walley predictive intervals &
Explicit selector: Bayes-normal, Bayes-exact, or Walley-IDM \\
\hline
Intensional: $\text{InhInt} = \alpha \cdot \text{Ext} + (1-\alpha) \cdot \text{Int}$ &
Binary mixed policy insufficient &
Typed ASSOC/PAT channels with separate endpoints \\
\hline
Multideduction: two-term IE suffices &
Identifiability no-go &
Residual decomposition: $|A \cup B| = T_2 + R$ \\
\hline
Trail-free: converges &
2-cycle counterexample &
Damped variant with fresh evidence \\
\hline
Pooling axioms: unique &
Max-pooling satisfies axioms too &
External Bayes + total additivity $\Rightarrow$ unique \\
\hline
\end{longtable}


\chapter{Counterexample Catalogue}
\label{app:counterexamples}

\begin{longtable}{p{0.30\textwidth}p{0.25\textwidth}p{0.35\textwidth}}
\textbf{Counterexample} & \textbf{Module} & \textbf{What it blocks} \\
\hline
\code{same\_first\_two\_terms\_different\_union} &
PLNInclusionExclusionIdentifiability &
Two-term IE identifying union mass \\
\hline
\code{orbit\_not\_eventually\_constant} &
PLNTrailFreeDynamicsCounterexample &
Trail-free stabilization without damping \\
\hline
\code{maxPoolingOperator\_ne\_fuse} &
PremiseSelectionExternalBayesianity &
Pooling axioms characterizing additive fusion \\
\hline
\code{sameCoverageAndCard\_differentFalsePositiveRisk} &
PremiseSelectionCoverageCounterexamples &
Coverage/cardinality surrogate for risk \\
\hline
\code{threshold\_or\_counterexample} &
OSLFEvidenceSemantics &
Disjunction threshold without total order \\
\hline
\code{threshold\_dia\_fails} &
OSLFEvidenceSemantics &
Diamond threshold without finite branching \\
\hline
\code{no\_binary\_mixed\_policy\_for\_assocPat\_fixture} &
PLNIntensionalCanary &
Collapsing ASSOC+PAT to binary law \\
\hline
No-go (local completeness) &
PLNJointEvidenceNoGo &
Computing exact link evidence from local premises \\
\hline
\end{longtable}


\chapter{Deferred Claims}
\label{app:deferred}

\begin{longtable}{p{0.25\textwidth}p{0.35\textwidth}p{0.30\textwidth}}
\textbf{Claim} & \textbf{Why deferred} & \textbf{Status} \\
\hline
Executable PLN engine &
Requires MeTTa/Hyperon runtime &
\verdictHeuristic{} \\
\hline
Arbitrary-domain quantifier parity &
Basic infinitary semantics exist; full theorem/canary parity remains &
Partial \\
\hline
Chapter~12 asserted claim set &
No longer deferred: claim audit complete; chapter formulas are now explicitly proved, strengthened, or ruled out &
Settled \\
\hline
Universal Prediction core &
21-file WIP module &
Not in stable narrative \\
\hline
Chapter~9 operational closure &
Theorem-level complete; runtime pending &
\verdictHeuristic{} \\
\hline
\end{longtable}


\chapter{Notation Translation}
\label{app:notation-translation}

\begin{longtable}{p{0.25\textwidth}p{0.30\textwidth}p{0.35\textwidth}}
\textbf{PLN book} & \textbf{Lean} & \textbf{\xiPLN{} notation} \\
\hline
$\langle s, c \rangle$ & \code{STV s c} & $\STV{s}{c}$ \\
Strength $s$ & \code{queryStrength} & $\strength(W, q)$ \\
Confidence $c$ & \code{queryConfidence} & $\confidence(W, q)$ \\
Evidence $(n^+, n^-)$ & \code{Evidence.mk nPos nNeg} & $\Evidence$ \\
Revision & \code{Evidence.add} & $e_1 \hplus e_2$ \\
$P(B | A)$ & \code{link A B} query & $\code{link}(A, B)$ \\
Inheritance $A \Rightarrow B$ & concept-side of \code{link} & Link query \\
Implication & predicate-side of \code{link} & Link query \\
World model & \code{WorldModel State Query} & $\WorldModel$ \\
Side conditions & \code{WorldModelSigma.sigma} & $\Sigma$ \\
Indefinite TV $[L, U]$ & \code{ITVSelector} choice & Typed ITV with selector \\
\hline
\end{longtable}


\bibliographystyle{plain}
\bibliography{references}

\end{document}
